Let \(P\in\mathcal P_n(\Lambda_n)\).  The triangular-union definition
places \(P\) in \(\mathcal P_{\kappa}^{\mathrm{base}}\) for at least one
\(\kappa\in(2,\overline\kappa]\), so it satisfies every member condition
of that base class.  Fix \(x\in\mathcal B\), \(t\in\{0,1\}\), and
\(s>0\), and put \(b=\min\{h_0,s/2\}\).  Polynomial lower mass gives
\[
 P_X\{z\in\mathcal A_t:d(z,x)<s\}
 \ge q_{t,P}(x,b)\ge c_mb^\kappa>0.
 \tag{95}
\]
Thus \(x\in\mathcal X_{t,P}\).  The Lipschitz member condition makes
\(g_{t,P}\) continuous at \(x\) on that support, so the trace definition
gives
\[
 \theta_{t,P}(x)=g_{t,P}(x).
 \tag{96}
\]

Write \(z_n=\sqrt{2\log(2M_n/\gamma)}\).  For fixed scores let
\(\mathscr S_n\) be the disk-pattern family in
\(\mathrm{lem:disk\text{-}pattern\text{-}sub\text{-}gaussian\text{-}control}\).
Its cardinality is at most \(M_n\).  Substitution of \(z_n\) into that
lemma gives, for almost every score vector,
\[
 \Pr_P\!\left\{
 \max_{S\in\mathscr S_n}
 \frac{|\sum_{i\in S}\{Y_i-g_{T_i,P}(X_i)\}|}
      {\sigma\sqrt{|S|}}>z_n
 \ \middle|\ X_1,\ldots,X_n\right\}
 \le2|\mathscr S_n|e^{-z_n^2/2}\le\gamma.
 \tag{97}
\]
Independent sampling and the restricted-side conditional sub-Gaussian
kernel condition are exactly the premises used by that lemma.
Complementarity ensures that a side-\(t\) pattern contains precisely
observations having \(T_i=t\).

For a selected side fit abbreviate
\[
 \widehat h=\widehat h_{t,n}(x),\qquad
 N=N_{t,n}(x,\widehat h),\qquad
 S=\{i:X_i\in\mathcal A_t, d(X_i,x)\le\widehat h\}.
\]
If \(N>0\), then \(S\in\mathscr S_n\), and outside the event in (97),
\[
 \left|\frac1N\sum_{i\in S}\{Y_i-g_{t,P}(X_i)\}\right|
 \le\sigma\sqrt{\frac{2\log(2M_n/\gamma)}N}.
 \tag{98}
\]
Equation (96) and Lipschitzness also give
\[
 \left|\frac1N\sum_{i\in S}g_{t,P}(X_i)-\theta_{t,P}(x)\right|
 \le\frac1N\sum_{i\in S}Ld(X_i,x)
 \le L\widehat h.
 \tag{99}
\]
Thus the side error is at most \(B_{t,n}(x)\).  If \(N=0\), the
local-constant fit is zero and bounded traces give
\[
 |\widehat g_{t,n}(x;\widehat h)-\theta_{t,P}(x)|
 \le L=B_{t,n}(x).
 \tag{100}
\]
The event in (97) controls all disk patterns simultaneously.  Hence
(98)--(100) hold for every \(x\) and both sides with conditional
probability at least \(1-\gamma\).  The estimator definition and
\(\tau_P=\theta_{1,P}-\theta_{0,P}\) yield
\[
 |\widehat\tau_n(x)-\tau_P(x)|
 \le B_{0,n}(x)+B_{1,n}(x),
 \qquad x\in\mathcal B.
 \tag{101}
\]
Integrating over the scores proves
\[
 \inf_{P\in\mathcal P_n(\Lambda_n)}
 P\{\tau_P(x)\in\mathcal C_n(x)\text{ for every }x\in\mathcal B\}
 \ge1-\gamma.
 \tag{102}
\]
The derivation of (95)--(101) used only the member properties of a single
base class, so the same argument also proves coverage at least
\(1-\gamma\) uniformly over each fixed
\(\mathcal P_\kappa^{\mathrm{base}}\), and hence over its pervasive
subclass.
Each count and numerator is a finite sum of Borel functions; the grid is
finite; and the bandwidth definition supplies a deterministic tie break
and empty-set value.  Therefore the bandwidths, fits, radii, and band
endpoints are data measurable.

For expected width fix \(\kappa\), put
\[
 r_n=(\log n/n)^{1/(\kappa+2)},
 \tag{103}
\]
and choose \(h_*\in\mathcal H_n\) with \(r_n\le h_*<2r_n\), which is
possible for all large \(n\).  A maximal \(h_*/2\)-separated boundary
set \(\mathcal N_*\) is an \(h_*/2\)-net and satisfies
\[
 |\mathcal N_*|\le2C_{\mathcal B}h_*^{-1}.
 \tag{104}
\]
For \(z\in\mathcal N_*\), define
\[
 \widetilde N_{t,n}(z)=\sum_{i=1}^n
 \mathbf1\{X_i\in\mathcal A_t, d(X_i,z)\le h_*/2\}.
\]
It is binomial under independent sampling, with success probability
\[
 q_{t,P}(z,h_*/2)\ge c_m(h_*/2)^\kappa.
 \tag{105}
\]
For \(Z\sim\operatorname{Bin}(n,q)\), Markov's inequality applied to
\(2^{-Z}\), the binomial generating function, and
\(\log(1-u)\le-u\) give
\[
 \Pr\{Z\le nq/2\}
 \le2^{nq/2}(1-q/2)^n
 \le e^{-nq/8}.
 \tag{106}
\]
Union-bound (106) over (104) and both sides.  If \(x\in\mathcal B\),
choose \(z\in\mathcal N_*\) within \(h_*/2\); then
\(B(z,h_*/2)\subseteq B(x,h_*)\).  We obtain an event \(G_n\) such that,
uniformly over \(P\in\mathcal P_\kappa^{\mathrm{base}}\),
\[
 \inf_{x\in\mathcal B,\ t\in\{0,1\}}N_{t,n}(x,h_*)
 \ge cnh_*^\kappa,
 \qquad
 P(G_n^c)\le Ch_*^{-1}e^{-cnh_*^\kappa}.
 \tag{107}
\]
From (103) and \(h_*\asymp r_n\),
\[
 nh_*^{\kappa+2}\asymp\log n,
 \qquad
 \frac{nh_*^\kappa}{\log(e+n/h_*)}\longrightarrow\infty.
 \tag{108}
\]
Thus \(h_*\) is admissible for every \(x,t\) on \(G_n\), eventually.
Minimization of the bandwidth criterion and comparison with \(h_*\)
give
\[
 L\widehat h_{t,n}(x)+
 \sigma\sqrt{\frac{\log(e+n/\widehat h_{t,n}(x))}
 {N_{t,n}(x,\widehat h_{t,n}(x))}}
 \le C r_n
 \quad\text{on }G_n.
 \tag{109}
\]
Indeed, the criterion at \(h_*\) is at most
\[
 Lh_*+\sigma\sqrt{\frac{\log(e+n/h_*)}{cnh_*^\kappa}}
 \le Cr_n.
 \tag{110}
\]

The elementary bounds \(|\mathcal H_n|\le2+\log_2n\) and
\(\sum_{j=0}^3\binom nj\le Cn^3\) yield
\[
 \log(2M_n/\gamma)\le C_\gamma\log n
 \le C_\gamma\log(e+n/\widehat h_{t,n}(x)),
 \tag{111}
\]
because \(0<\widehat h_{t,n}(x)\le h_0<1\).  On \(G_n\) the selected
count is positive, so (109)--(111) imply
\[
 \sup_{x\in\mathcal B,\ t\in\{0,1\}}B_{t,n}(x)
 \le C_\gamma r_n.
 \tag{112}
\]
On \(G_n^c\), the two branches in the definition of \(B_{t,n}\),
\(h_0<1\), and \(N\ge1\) in the positive-count branch give
\[
 \sup_{x,t}B_{t,n}(x)
 \le L+\sigma\sqrt{2\log(2M_n/\gamma)}.
 \tag{113}
\]
Moreover,
\[
 h_*^{-1}e^{-cnh_*^\kappa}
 \sqrt{\log(2M_n/\gamma)}=o(r_n),
 \tag{114}
\]
because \(nh_*^\kappa\) grows as a positive power of \(n\), up to a
logarithmic factor.  Since the displayed interval has
\[
 W_n(x)=B_{0,n}(x)+B_{1,n}(x),
 \tag{115}
\]
splitting the expectation over \(G_n\) and its complement proves
\[
 \sup_{P\in\mathcal P_\kappa^{\mathrm{base}}}
 \mathbb E_P\sup_{x\in\mathcal B}W_n(x)
 \le C_\gamma(\log n/n)^{1/(\kappa+2)}.
 \tag{116}
\]

Suppose now that \(\mathcal P_\kappa^{\mathrm{perv}}\ne\varnothing\),
and let a band have simultaneous coverage at least
\(1-\gamma-\epsilon_n\) on that class, with \(\epsilon_n\downarrow0\).
The assignment regions are complementary, the geometry satisfies global
boundary entropy, and the class supplies all base conditions and
pervasive thin sites.  Therefore
\(\mathrm{lem:geometry\text{-}general\text{-}pervasive\text{-}band\text{-}width\text{-}obstruction}\)
gives
\[
 \sup_{P\in\mathcal P_\kappa^{\mathrm{perv}}}
 \mathbb E_P\sup_{x\in\mathcal B}W_n(x)
 \ge c_\gamma(\log n/n)^{1/(\kappa+2)}
 \tag{117}
\]
and
\[
 \sup_{P\in\mathcal P_\kappa^{\mathrm{perv}}}
 \sup_{x\in\mathcal B}
 \frac{\mathbb E_PW_n(x)}{r_{n,P}^{\mathrm{loc}}(x)}
 \ge c_\gamma.
 \tag{118}
\]
These conclusions are geometry-general; the circular membership lemma
only certifies explicit nonemptiness.  Since
\(\mathcal P_\kappa^{\mathrm{perv}}\subseteq
\mathcal P_\kappa^{\mathrm{base}}\), honesty over the base class implies
honesty over the subclass and the base-class supremum dominates (117).
Conversely, (116) applies to the explicit band on both classes.  This
proves the matched global expected-supremum half-width frontier, while
(118) is the claimed local-benchmark converse.

It remains to prove the general two-law inequality.  Fix
\(x\in\mathcal B\), and for \(j=0,1\) let
\(A_j=\{\tau_{P_j}(x)\in\mathcal C_n(x)\}\).  Simultaneous coverage over
a class containing both laws implies
\[
 P_j^{\otimes n}(A_j)\ge1-\gamma-\epsilon_n,
 \qquad j=0,1.
 \tag{119}
\]
The defining variational inequality for total variation of the full
observed product laws gives
\[
 P_0^{\otimes n}(A_1)
 \ge P_1^{\otimes n}(A_1)
 -\operatorname{TV}(P_0^{\otimes n},P_1^{\otimes n}).
 \tag{120}
\]
Together with the union bound in reverse, (119)--(120) yield
\[
 P_0^{\otimes n}(A_0\cap A_1)
 \ge1-2\gamma-2\epsilon_n
 -\operatorname{TV}(P_0^{\otimes n},P_1^{\otimes n}).
 \tag{121}
\]
On \(A_0\cap A_1\), the interval contains both target values, so
\[
 W_n(x)\ge\frac12|\tau_{P_1}(x)-\tau_{P_0}(x)|.
 \tag{122}
\]
Multiply (122) by the intersection indicator, take the
\(P_0\)-expectation, and use the positive part of (121).  This proves
\[
 \mathbb E_{P_0}W_n(x)\ge
 \frac{|\tau_{P_1}(x)-\tau_{P_0}(x)|}{2}
 \left[1-2\gamma-2\epsilon_n-
 \operatorname{TV}(P_0^{\otimes n},P_1^{\otimes n})\right]_+.
\]
The argument necessarily compares full observed laws.  No step converts
score-law total variation into full-law distinguishability, and the
global logarithm in (111) need not equal the bounded local-multiplicity
term in \(r_{n,P}^{\mathrm{loc}}\).  Thus the proof establishes exactly
the global and obstruction claims, without asserting the stated open
local-oracle or score-law-only conclusions.